\documentclass[11pt]{amsart}
\usepackage[T1]{fontenc}
\usepackage{lmodern}
\usepackage{amsmath,amssymb,amsthm}
\usepackage[numbers,sort]{natbib}
\usepackage[hidelinks]{hyperref}
\usepackage{bm}

\newtheorem{theorem}{Theorem}
\newtheorem{lemma}[theorem]{Lemma}
\newtheorem{proposition}[theorem]{Proposition}
\newtheorem{corollary}[theorem]{Corollary}
\theoremstyle{definition}
\newtheorem{definition}[theorem]{Definition}
\newtheorem{remark}[theorem]{Remark}

\newcommand{\Ter}{\mathcal S}
\newcommand{\red}{\twoheadrightarrow}
\newcommand{\hred}{\twoheadrightarrow_h}
\newcommand{\conv}{=_\beta}
\newcommand{\nconv}{\neq_\beta}
\newcommand{\sz}[1]{\lvert#1\rvert}
\newcommand{\NF}{\mathrm{NF}}
\newcommand{\rest}{\mathrm{rest}}
\newcommand{\B}{\mathbb B}
\newcommand{\Hd}{\mathrm{H}}
\newcommand{\dom}{\mathrm{dom}}
\newcommand{\Bad}{\mathrm{B}}
\newcommand{\li}{\mathsf i}
\newcommand{\la}{\mathsf a}
\newcommand{\lb}{\mathsf b}
\newcommand{\lr}{\mathsf r}
\newcommand{\lw}{\mathsf w}
\newcommand{\CT}{\mathrm{CT}}

\title[Spines and Turns]{Convertibility of Tree-Equal $\bm{S}$-Terms Splits into Spines and Turns}
\author{Claude Opus \& Fable with N.\ E.\ Davis}
\address{ZeroAttest, Champaign, Illinois, USA}
\email{NDavis@ZeroAttest.com}
\urladdr{https://orcid.org/0000-0002-1131-8319}
\date{\today}

\begin{document}
\begin{abstract}
Two closed terms built from the combinator $S$ are tree-equal when their infinitary normal forms
agree, and Lock~1 is the conjecture that tree-equal terms are convertible. We split Lock~1 into two
conjectures of different shape. A term without a head normal form splits off an infinite sequence
of arguments, its teeth, along an infinite left spine; CT says that two such terms with pairwise
convertible teeth are convertible. Two tree-equal terms that are not convertible have a pair of
children that again are tree-equal and not convertible, so they have an infinite descent through
such pairs; NTD says that no such descent leaves a left spine infinitely often. We prove that Lock~1 holds if and
only if CT and NTD both hold. The forward direction is constructive; the converse uses excluded
middle for convertibility. Three short results locate CT. CT at a term is equivalent to CT at the
next term of its spine, with no further hypothesis, and, when Lock~1 holds on the classes of its
teeth, to Lock~1 on its own class. Lock~1 on the class of a term passes to every node of its tree;
so Lock~1 on the class of $D\,(c\,D)$, where $c=S\,(S\,S)$ and $D=S\,c\,c$, implies CT at infinitely
many explicit terms, the smallest of sizes $50$ and $56$, and a counterexample to CT at one of them
gives an explicit pair of tree-equal terms that are not convertible. A census shows how little of
CT is known: of the $65{,}794$ terms of sizes $8$ to $13$ without a head normal form, a computation
certifies only $535$ as covered by the known cases of CT. Lock~1, CT and NTD are open. Nothing in
this paper is mechanized; the congruence theorem for tree equality that it uses is.
\end{abstract}

\keywords{combinatory logic, combinator $S$, infinitary normal forms, convertibility, head normal forms}
\maketitle

\section{Introduction}

Let $\Ter$ be the set of closed terms built from $S$, with $S\,x\,y\,z\to x\,z\,(y\,z)$ at any
position, and write $\conv$ for convertibility. Every $t\in\Ter$ has an infinitary normal form
$t^\infty$~\citep{s3}, and we write $t\approx u$, \emph{$t$ and $u$ are tree-equal}, if $t$ and $u$
have the same finite prefixes $\NF_d$ at every depth~\citep[Definition~1]{s13}. Convertible terms are
tree-equal. The converse is
\begin{center}
\emph{Lock~1:}\enskip $t\approx u$ implies $t\conv u$,
\end{center}
the conjecture of~\citep[Conjecture~11]{s6}. It would make convertibility of $S$-terms
decidable~\citep{s6}, and it is open.

A term $p$ without a head normal form (hnf) has an infinite left spine: its head run reaches a first
stable frame $\rest(p)\,\tau_1(p)$, the term $\rest(p)$ again has no hnf, and so on, and the
arguments split off on the way are the \emph{teeth} $\tau_n(p)$ of $p$~\citep{s11}. The tree of $p$ is
the infinite left comb of the trees of its teeth (Lemma~\ref{lem:bot}).

\begin{definition}[CT~\citep{s11,s12}]\label{def:CT}
For $R\in\Ter$ without an hnf, $\CT(R)$ says: every $u\in\Ter$ without an hnf such that
$\tau_n(u)\conv\tau_n(R)$ for all $n\ge1$ satisfies $u\conv R$. \emph{CT} says that $\CT(R)$ holds for
every such $R$.
\end{definition}

A pair $t\approx u$ with $t\nconv u$ is \emph{bad}. Every bad pair has a bad pair among its
children (Lemma~\ref{lem:bad}); so it has an infinite descent through bad pairs. A step
of a descent that is not a step down a left spine is a \emph{turn}
(Definition~\ref{def:descent}).

\begin{definition}[NTD]\label{def:NTD}
\emph{No turning descent:} for no $t\approx u$ does the bad set $\Bad(t,u)$ contain a descent with
infinitely many turns.
\end{definition}

\begin{theorem}\label{thm:main}
Lock~1 holds if and only if CT and NTD hold.
\end{theorem}

A counterexample to CT is a bad pair whose only descent runs down one left spine with no turn
(Lemma~\ref{lem:spine}); NTD forbids exactly the descents that turn infinitely often. So the theorem
splits the possible counterexamples to Lock~1 by the shape of their descents. It is a decomposition
into two consequences of Lock~1, not a reduction to a known statement: all three statements are
open.

Section~\ref{sec:ct} gives two short facts about CT at one term $R$. $\CT(R)$ is equivalent to
$\CT(\rest(R))$ (Proposition~\ref{prop:shift}). If Lock~1 holds on the $\approx$-class of every tooth of
$R$, then $\CT(R)$ is equivalent to Lock~1 on the $\approx$-class of $R$
(Proposition~\ref{prop:bot}). Section~\ref{sec:her} shows that Lock~1 on the $\approx$-class of a term
passes to every node of its tree (Lemma~\ref{lem:replace}) and applies this to a family of terms
whose trees meet nodes without an hnf two steps off a path of nodes with an hnf. Section~\ref{sec:known}
records how few cases of CT are known, and Section~\ref{sec:scope} what is not proved.

\emph{Assumptions.} We use the facts on blocks and prefixes listed in~\citep[\S2]{s13}; Fact~K and
Lemma~2 of~\citep{s13}; Theorem~C of~\citep{s13}, that $\approx$ is a congruence; and
Lemmas~3--6 of~\citep{s11} on head runs, teeth and convertibility. All of these are mechanized in
the cited papers. Section~\ref{sec:known} cites the CT theorems of~\citep{s11,s12}; the proofs
of~\citep{s12} are not mechanized. The converse direction of Theorem~\ref{thm:main} uses excluded
middle for $\conv$, at the steps marked; all other proofs are constructive, since every case split
on whether a term has an hnf is the decidable test of its block~\citep{s1}.

\section{Preliminaries}\label{sec:prelim}

Every $t\in\Ter$ is $S\,u_1\cdots u_m$ for a unique $m\ge0$; the $u_i$ are its \emph{items}, and $t$
is an \emph{hnf} if $m\le2$. If $m\ge3$ the \emph{head step} is $t\to u_1\,u_3\,(u_2\,u_3)\,u_4\cdots
u_m$; $\hred$ is the reflexive--transitive closure of head steps, and the terms of the head run are
its \emph{frames}. A term \emph{has an hnf} if its head run reaches one; this holds iff its block in the
38-element model $\B$ of~\citep{s1} is not $\bot=36$, and $\beta(\bot,x)=\bot$ for every $x$. The
\emph{size} $\sz t$ is the number of occurrences of $S$ in $t$.

\emph{Teeth}~\citep[\S1]{s11}. A frame $f\,w$ is \emph{stable} if $f$ has no hnf. Let $t$ have no hnf.
Its head run reaches a stable frame; if $f\,w$ is the first, then $\rest(t)=f$ and $\tau_1(t)=w$, and
$\rest(t)$ has no hnf. Put $\tau_{n+1}(t)=\tau_1(\rest^n(t))$. We call $t$ \emph{literal} if it is its
own first stable frame, that is, $t=f\,w$ with $f$ without an hnf, and \emph{reduced} otherwise.
From~\citep{s11} we use:
\begin{itemize}
\item[(L1)] if $f$ has no hnf, then $f\,w$ has none, and $f\,w$ is literal with $\rest(f\,w)=f$ and
  $\tau_1(f\,w)=w$ \citep[Lemmas~3(d) and~4(A)]{s11};
\item[(L2)] if $t\conv t'$ and $t$ has the hnf $S\,u_1\cdots u_m$, then $t'$ has an hnf $S\,u'_1\cdots
  u'_m$ with $u_i\conv u'_i$; if $t\conv t'$ and $t$ has no hnf, then $t'$ has none, and
  $\rest^n(t)\conv\rest^n(t')$ and $\tau_n(t)\conv\tau_n(t')$ for all $n$~\citep[Lemma~5]{s11};
\item[(L3)] if $t$ has no hnf, then $t\hred\rest^n(t)\,\tau_n(t)\cdots\tau_1(t)$; hence, for
  $p,q$ without an hnf, $\rest^n(p)\conv\rest^n(q)$ and $\tau_i(p)\conv\tau_i(q)$ for all
  $i\le n$ imply $p\conv q$~\citep[Lemma~6 and its proof]{s11}.
\end{itemize}

\emph{Trees}~\citep[\S2]{s13}. $\NF_d(t)$ is the prefix of depth $d$; $t\approx u$ iff
$\NF_d(t)=\NF_d(u)$ for all $d$. We use: (N2) if $t\red u$ then $\NF_d(t)=\NF_d(u)$, so $\conv$
implies $\approx$; (N4) if $t\hred f\,w$ and $f\,w$ is stable or an hnf, then
$\NF_{d+1}(t)=\NF_d(f)\cdot\NF_d(w)$; and
\begin{itemize}
\item[(K1)] if $t\approx t'$, then $t$ has an hnf iff $t'$ has one~\citep[Lemma~2(c)]{s13}; if $t\hred
  f\,w$ and $t'\hred f'\,w'$ are stable, then $f\approx f'$ and $w\approx
  w'$~\citep[Lemma~2(b)]{s13};
\item[(K2)] \emph{Fact K}: if $t\approx t'$ and $t$ has the hnf $S\,u_1\cdots u_m$, then $t'$ has an
  hnf $S\,u'_1\cdots u'_m$ with $u_i\approx u'_i$~\citep[Lemma~3]{s13};
\item[(C)] \emph{Theorem~C}: if $t\approx t'$ and $u\approx u'$, then $t\,u\approx
  t'\,u'$~\citep{s13}.
\end{itemize}
We say that \emph{Lock~1 holds on the class of $t$} if every $u\approx t$ satisfies $u\conv t$.

\begin{lemma}\label{lem:bot}
Let $R$ have no hnf. Then $u\approx R$ iff $u$ has no hnf and $\tau_n(u)\approx\tau_n(R)$ for all
$n\ge1$. If $u\approx R$, then also $\rest^n(u)\approx\rest^n(R)$ for all $n$.
\end{lemma}

\begin{proof}
If $u\approx R$, then $u$ has no hnf by (K1), and the first stable frames give $\rest(u)\approx\rest(R)$
and $\tau_1(u)\approx\tau_1(R)$ by (K1); iterate. Conversely, we show by induction on $d$ that
$\NF_d(\rest^n(u))=\NF_d(\rest^n(R))$ for all $n$. For $d=0$ both are a hole. By (N4) at the first
stable frame of $\rest^n(u)$,
$\NF_{d+1}(\rest^n(u))=\NF_d(\rest^{n+1}(u))\cdot\NF_d(\tau_{n+1}(u))$, and likewise for $R$; the
induction hypothesis and $\tau_{n+1}(u)\approx\tau_{n+1}(R)$ make the two equal. Take $n=0$.
\end{proof}

This proves the equivalence sketched in~\citep[Remark~10]{s14}.

\section{The descent}\label{sec:descent}

\begin{definition}\label{def:descent}
The \emph{kind} of $t$ is $m\in\{0,1,2\}$ if $t$ has the hnf $S\,u_1\cdots u_m$, and $\bot$ if $t$ has
no hnf. The \emph{head form} $\Hd(t)$ is that hnf, or the first stable frame. The \emph{children} of
$t$ are: none at kind $0$; $t_\li=a$ if $\Hd(t)=S\,a$; $t_\la=a$ and $t_\lb=b$ if $\Hd(t)=S\,a\,b$;
$t_\lr=\rest(t)$ and $t_\lw=\tau_1(t)$ at kind $\bot$. For a word $\sigma$ over
$\{\li,\la,\lb,\lr,\lw\}$, $t_\sigma$ is defined by $t_\varepsilon=t$ and $t_{\sigma x}=(t_\sigma)_x$, and
$\dom(t)$ is the set of $\sigma$ for which it is defined. A \emph{turn} is a letter other than $\lr$.
For $t\approx u$ the \emph{bad set} is $\Bad(t,u)=\{\sigma\in\dom(t):t_\sigma\nconv u_\sigma\}$. A
\emph{descent} in a set $W$ of words is an infinite word all of whose finite prefixes lie in $W$.
\end{definition}

In the tree of $t$, the letters $\lr$ and $\lw$ go to the left and right subtree, $\li$ and $\lb$ to
the right subtree, and $\la$ to the right subtree of the left subtree. So $t_{\lr^{n-1}\lw}=\tau_n(t)$.
Since $t\hred\Hd(t)$, we have $t\conv\Hd(t)$ and $t\approx\Hd(t)$.

\begin{lemma}\label{lem:kids}
\textup{(a)} If $t\approx u$, then $t$ and $u$ have the same kind and corresponding children are
tree-equal; hence $\dom(t)=\dom(u)$ and $t_\sigma\approx u_\sigma$ for all $\sigma\in\dom(t)$.

\textup{(b)} If $t\conv u$, then $t$ and $u$ have the same kind and corresponding children are
convertible.

\textup{(c)} If $t$ and $u$ have the same kind and corresponding children are convertible, then
$t\conv u$.
\end{lemma}

\begin{proof}
(a) (K1) and (K2); induct on $\sigma$. (b) (L2). (c) $t\conv\Hd(t)\conv\Hd(u)\conv u$, since
$\conv$ is a congruence.
\end{proof}

\begin{lemma}\label{lem:bad}
Let $t\approx u$ and $\Bad=\Bad(t,u)$.

\textup{(a)} $\Bad$ is closed under prefixes.

\textup{(b)} No $\sigma$ with $t_\sigma$ of kind $0$ lies in $\Bad$. If $\sigma\in\Bad$, then
$\sigma x\in\Bad$ for some letter $x$ (excluded middle for $\conv$).
\end{lemma}

\begin{proof}
(a) If $t_\sigma\conv u_\sigma$, then $t_{\sigma\rho}\conv u_{\sigma\rho}$ for all $\rho$ by
Lemma~\ref{lem:kids}(b). (b) At kind $0$ both terms reduce to $S$. If $\sigma\in\Bad$, the pairs of
children are not all convertible, by Lemma~\ref{lem:kids}(a,c); excluded middle for each of the at
most two pairs gives one that is not convertible.
\end{proof}

\section{Spines and the proof of Theorem~\ref{thm:main}}\label{sec:proof}

A \emph{spine pair} is a pair $p\nconv q$ of terms without an hnf with $\tau_n(p)\conv\tau_n(q)$ for all
$n\ge1$: a counterexample to CT.

\begin{lemma}\label{lem:spine}
A spine pair $(p,q)$ satisfies $p\approx q$, and $\Bad(p,q)=\{\lr^n:n\ge0\}$.
\end{lemma}

\begin{proof}
$p\approx q$ by Lemma~\ref{lem:bot}, since $\conv$ implies $\approx$. If $\rest^n(p)\conv\rest^n(q)$,
then $p\conv q$ by (L3); so every $\lr^n$ is bad. The teeth are convertible, so no word with a letter
$\lw$ is bad, by Lemma~\ref{lem:bad}(a).
\end{proof}

\begin{proof}[Proof of Theorem~\ref{thm:main}]
\emph{Lock~1 implies CT.} If $p,q$ have no hnf and convertible teeth, then $p\approx q$ by
Lemma~\ref{lem:bot}, and $p\conv q$ by Lock~1. \emph{Lock~1 implies NTD.} If $t\approx u$, then
$t_\sigma\approx u_\sigma$ for every $\sigma\in\dom(t)$ (Lemma~\ref{lem:kids}(a)), so Lock~1 gives
$\Bad(t,u)=\emptyset$, which contains no descent. Both arguments are constructive.

\emph{CT and NTD imply Lock~1.} Let $t\approx u$, and suppose $t\nconv u$. To every bad pair $(x,y)$ we
assign a nonempty word $\delta\in\Bad(x,y)$ that ends with a turn. By Lemma~\ref{lem:bad}(b), $x$ is not
of kind $0$. At kind $1$ or $2$, let $\delta$ be the first letter in the order $\li,\la,\lb$ with
$\delta\in\Bad(x,y)$; it exists by Lemma~\ref{lem:bad}(b), and finding it uses excluded middle for
$\conv$. At kind $\bot$, $(x,y)$ is not a spine pair, by $\CT(y)$; so it is not the case that
$\tau_n(x)\conv\tau_n(y)$ for every $n$, and excluded middle for $\conv$ gives a least $n$ with
$\tau_n(x)\nconv\tau_n(y)$. Let $\delta=\lr^{n-1}\lw$. In both cases $\delta\in\Bad(x,y)$, and
$(x_\delta,y_\delta)$ is again bad, by Lemma~\ref{lem:kids}(a).

Put $\sigma_0=\varepsilon$ and $\sigma_{k+1}=\sigma_k\,\delta(t_{\sigma_k},u_{\sigma_k})$. Each
$\sigma_k$ lies in $\Bad(t,u)$, since $(t_{\sigma_k})_\rho=t_{\sigma_k\rho}$. The infinite word whose
prefixes are the $\sigma_k$ is a descent in $\Bad(t,u)$, by Lemma~\ref{lem:bad}(a), and it has
infinitely many turns, one at the end of each $\delta$. This contradicts NTD. So $t\nconv u$ is
refuted. The choices of $\delta$ are definite (a first letter, a least index), so no choice principle
beyond excluded middle for $\conv$ is used.

From $\neg\neg(t\conv u)$ we conclude $t\conv u$. This last step needs only Markov's principle,
because $t\conv u$ says that a search for a common reduct succeeds; the construction of the descent
before it uses excluded middle.
\end{proof}

Read contrapositively, and classically: if Lock~1 fails, then a spine pair exists, or some bad pair
has a descent with infinitely many turns. NTD says nothing about a spine pair, whose only descent has
no turn (Lemma~\ref{lem:spine}); whether NTD implies CT is open.

\section{CT at one term}\label{sec:ct}

\begin{proposition}\label{prop:shift}
Let $R$ and $R'$ have no hnf.

\textup{(a)} If $R\conv R'$, then $\CT(R)$ iff $\CT(R')$.

\textup{(b)} $\CT(R)$ iff $\CT(\rest(R))$.

Hence, if $\rest^i(R)\conv\rest^j(R')$ for some $i,j$, then $\CT(R)$ iff $\CT(R')$.
\end{proposition}

\begin{proof}
(a) The teeth of $R$ and $R'$ are convertible, and so are $R$ and $R'$ (L2); use transitivity.
(b) Assume $\CT(R)$, and let $v$ have no hnf with $\tau_n(v)\conv\tau_{n+1}(R)$ for all $n$. Put
$u=v\,\tau_1(R)$. By (L1), $u$ has no hnf, $\rest(u)=v$ and $\tau_1(u)=\tau_1(R)$, so
$\tau_{n+1}(u)=\tau_n(v)$, and all teeth of $u$ are convertible with those of $R$. So $u\conv R$, and
$v=\rest(u)\conv\rest(R)$ by (L2). Conversely, assume $\CT(\rest(R))$, and let $u$ have no hnf and the
teeth of $R$. The teeth of $\rest(u)$ are $\tau_{n+1}(u)\conv\tau_{n+1}(R)$, so $\rest(u)\conv\rest(R)$,
and $u\conv R$ by (L3) with $n=1$.
\end{proof}

So $\CT(R)$ depends only on the tail of the spine of $R$, up to convertibility. The proof is
constructive and assumes nothing about the teeth.

\begin{proposition}\label{prop:bot}
Let $R$ have no hnf. Consider:
\begin{itemize}
\item[(i)] $\CT(R)$;
\item[(ii)] Lock~1 holds on the class of $R$;
\item[(iii)] for every $n$, every reduced $P\approx\rest^n(R)$ satisfies $P\conv\rest^n(R)$.
\end{itemize}
Then \textup{(ii)} implies \textup{(i)} and \textup{(iii)}. If Lock~1 holds on the class of every
tooth $\tau_n(R)$, the three are equivalent.
\end{proposition}

\begin{proof}
(ii)$\Rightarrow$(i). If $u$ has no hnf and the teeth of $R$, then $u\approx R$ by Lemma~\ref{lem:bot}.

(ii)$\Rightarrow$(iii). Let $P\approx\rest^n(R)$, reduced or not; $P$ has no hnf by (K1). Write
$\vec\tau$ for $\tau_n(R)\cdots\tau_1(R)$. By (C), $P\,\vec\tau\approx\rest^n(R)\,\vec\tau$, and by
(L3) and (N2) this is $\approx R$. So $P\,\vec\tau\conv R$. By (L1), $\rest^n(P\,\vec\tau)=P$, and
(L2) gives $P\conv\rest^n(R)$.

Now assume Lock~1 on the class of every tooth. (i)$\Rightarrow$(ii): if $u\approx R$, then $u$ has no hnf
and $\tau_n(u)\approx\tau_n(R)$ (Lemma~\ref{lem:bot}), so $\tau_n(u)\conv\tau_n(R)$, and $\CT(R)$
gives $u\conv R$. (iii)$\Rightarrow$(ii): let $u\approx R$. If $\rest^n(u)$ is literal, then
$\rest^{n+1}(u)$ is a proper subterm of it; so for some $n_0\le\sz u$ the term $\rest^{n_0}(u)$ is
reduced, and the least such $n_0$ is found by testing blocks. By Lemma~\ref{lem:bot},
$\rest^{n_0}(u)\approx\rest^{n_0}(R)$ and $\tau_i(u)\approx\tau_i(R)$, hence $\tau_i(u)\conv\tau_i(R)$.
By (iii), $\rest^{n_0}(u)\conv\rest^{n_0}(R)$, and $u\conv R$ by (L3).
\end{proof}

All steps are constructive. Under the hypothesis on the teeth, any argument that settles Lock~1 on the
class of $R$, for instance a certificate for the node $R$ of a tree, proves $\CT(R)$; no rule for
nodes without an hnf is weaker than CT at that node. Statement~(iii) is needed at every $n$: we do not
know whether it suffices at $n=0$ alone. Every term without an hnf has infinitely many reduced
$\rest^n$, by the subterm argument above, and at a reduced level (iii) is never vacuous, since
$\rest^n(R)$ is then a reduced member of its own class.

\section{Heredity}\label{sec:her}

\begin{lemma}\label{lem:replace}
Let $\sigma\in\dom(t)$ and $v\approx t_\sigma$. There is $t'$ with $t'\approx t$, such that
$t'\conv t$ implies $v\conv t_\sigma$. Hence if Lock~1 holds on the class of $t$, it holds on the class
of every $t_\sigma$.
\end{lemma}

\begin{proof}
By induction on $\sigma$. For $\sigma=\varepsilon$ take $t'=v$. For $\sigma=x\rho$, let $c=t_x$, and let
$c'$ be given by the induction hypothesis for $c$, $\rho$ and $v$: $c'\approx c$, and $c'\conv c$
implies $v\conv t_\sigma$. Let $t'$ be $\Hd(t)$ with the child $c$ replaced by $c'$. By (C) and
$\Hd(t)\approx t$, $t'\approx t$. The term $t'$ is its own head form with $x$-child $c'$ and its other
children as in $\Hd(t)$: an hnf is its own hnf, and if $x\in\{\lr,\lw\}$, then $t'$ is literal by (L1),
because its function part, $\rest(t)$ or $c'\approx\rest(t)$, has no hnf (K1). If $t'\conv t$, then
$t'\conv\Hd(t)$, and Lemma~\ref{lem:kids}(b) gives $c'\conv c$, so $v\conv t_\sigma$. For the last
sentence, if $v\approx t_\sigma$, then $t'\approx t$, so $t'\conv t$, so $v\conv t_\sigma$.
\end{proof}

The construction is explicit, and the proof is constructive. It is the injectivity of the map
that sends the class of $t$ to the classes of its children, iterated along $\sigma$.

\begin{corollary}\label{cor:her}
Let $\sigma\in\dom(t)$, and let $R=t_\sigma$ have no hnf. If Lock~1 holds on the class of $t$, then
$\CT(R)$. Explicitly, by contraposition with a witness: if $u$ has no hnf, $\tau_n(u)\conv\tau_n(R)$ for
all $n$, and $u\nconv R$, then the term $t'$ of Lemma~\ref{lem:replace} for $v=u$ satisfies $t'\approx t$
and $t'\nconv t$, and it agrees with $\Hd(t)$ except in the child $t_x$, where $\sigma=x\rho$.
\end{corollary}

\begin{proof}
Lemma~\ref{lem:replace} and Proposition~\ref{prop:bot}, (ii)$\Rightarrow$(i), which needs no
hypothesis on the teeth. For the witness, $u\approx R$ by Lemma~\ref{lem:bot}, so Lemma~\ref{lem:replace}
applies to $v=u$; if $t'\conv t$, then $u\conv R$.
\end{proof}

The converse, that CT at the nodes without an hnf implies Lock~1 on the class of $t$, is not claimed.

\emph{An example.} Let $c=S\,(S\,S)$, $D=S\,c\,c$, $H_0=D$, $H_{j+1}=c\,H_j$, and $N_a(n,j)=H_j\,H_{n+1}$
for $0\le j\le n$. The trees of these terms have an infinite path of nodes with an hnf, and whether
Lock~1 holds on their classes is open~\citep[Remark~9]{s14}.

\begin{proposition}\label{prop:afam}
Let $0\le j\le n$. If $j\ge1$, put $A=N_a(n,j-1)$ and $m=n+1$; if $j=0$, put $A=N_a(n+1,n+1)$ and
$m=n+2$. Then $m\ge2$,
\[
N_a(n,j)\hred S\,A\,(H_m\,A),\qquad H_m\,A\hred S\,X\,R(n,j),
\]
where $X=H_{m-1}\,A$ and $R(n,j)=A\,(H_{m-1}\,A)$, and $R(n,j)$ has no hnf. So $R(n,j)=N_a(n,j)_{\lb\lb}$,
and $N_a(n,j)_\la=A$. Consequently:
\begin{itemize}
\item[(a)] if Lock~1 holds on the class of $N_a(n,j)$, then $\CT(R(n,j))$; and if $u$ has no hnf,
  $\tau_k(u)\conv\tau_k(R(n,j))$ for all $k$, and $u\nconv R(n,j)$, then $S\,A\,(S\,X\,u)\approx
  N_a(n,j)$ and $S\,A\,(S\,X\,u)\nconv N_a(n,j)$, with the same first item $A$;
\item[(b)] if Lock~1 holds on the class of $N_a(0,0)=D\,(c\,D)$, then $\CT(R(n,j))$ for all $0\le j\le n$.
\end{itemize}
The sizes are $\sz{R(n,j)}=35+9n+6j$ for $j\ge1$ and $56+15n$ for $j=0$; the two smallest are $50$, at
$(1,1)$, and $56$, at $(0,0)$.
\end{proposition}

\begin{proof}
$m\ge2$ because $j\ge1$ forces $n\ge1$. Two head steps give $c\,x\,y\hred S\,(x\,y)\,(y\,(x\,y))$.
For $j\ge1$, $N_a(n,j)=c\,H_{j-1}\,H_{n+1}\hred S\,A\,(H_{n+1}\,A)$. For $j=0$, $D\,W\to c\,W\,(c\,W)\hred
S\,(W\,(c\,W))\,(c\,W\,(W\,(c\,W)))$ with $W=H_{n+1}$, and $W\,(c\,W)=H_{n+1}\,H_{n+2}=A$. Then $H_m\,A=c\,H_{m-1}\,A\hred
S\,(H_{m-1}\,A)\,(A\,(H_{m-1}\,A))$. Blocks: $b(c)=2$, $b(D)=22$, and the table entries
$\beta(2,22)=\beta(2,30)=30$, $\beta(22,30)=\beta(30,30)=31$, $\beta(30,31)=31$ and $\beta(31,31)=36$ give,
by induction, $b(H_j)=30$ for $j\ge1$, $b(N_a(n,j))=31$ for all $n,j$, and, as $m-1\ge1$,
$b(R(n,j))=\beta(31,\beta(30,31))=\beta(31,31)=36$. So $R(n,j)$ has no hnf.

(a) is Corollary~\ref{cor:her} with $\sigma=\lb\lb$: the head forms $\Hd(N_a(n,j))=S\,A\,(H_m\,A)$ and
$\Hd(H_m\,A)=S\,X\,R(n,j)$ are replaced by $S\,A\,(S\,X\,u)$. (b) The first item of $N_a(n,j)$ is
$N_a(n,j-1)$ for $j\ge1$ and $N_a(n+1,n+1)$ for $j=0$, so $N_a(0,0)_{\la^k}$ runs through every
$N_a(n,j)$. Apply Lemma~\ref{lem:replace} with $\sigma=\la^k$, then (a). The sizes follow from
$\sz{H_j}=7+3j$.
\end{proof}

The proof uses six entries of the table $\beta$ of~\citep{s1}, read from it by a short program. We
also checked by direct reduction, for $n\le4$, the head forms, the identity of $R(n,j)$, its block,
and its size.

So any proof of Lock~1 on the one class of $D\,(c\,D)$, a term of size $17$, must prove CT at the
infinitely many terms $R(n,j)$. An argument that inspects only the path of first items cannot see
them: by~(a), a counterexample to CT at $R(n,j)$ gives a bad pair that agrees with $N_a(n,j)$ along
that whole path.

\section{How little of CT is known}\label{sec:known}

\begin{remark}[The known cases]\label{rem:known}
CT is proved at one term and at all terms whose run settles into an explicit grammar. Let $Q=S\,(S\,S)$,
$G=S\,Q\,Q$ and $p_K=S\,(S\,S)\,G\,(Q\,(Q\,G))\,S$. Then $\CT(R)$ holds when $R\conv\rest^n(p_K)$ for
some $n$~\citep[Theorem~1 and Corollary~2]{s11}, and when $\rest^e(R)\conv Z\,W$ for some $e$ and a pair
$(Z,W)$ of the grammar of~\citep[Definition~1]{s12} with all teeth of $Z\,W$
normalizing~\citep[Theorem~2]{s12}; \citep[Corollaries~3 and~4]{s12} give instances, among them the
runs of $X\,X$ with $X=S\,T\,S$ and $T=S\,S\,S$. By Proposition~\ref{prop:shift}, each case extends to
every $R$ with $\rest^i(R)\conv\rest^j(R')$ for such an $R'$.
\end{remark}

\begin{remark}[A census]\label{rem:census}
We enumerated the terms without an hnf of sizes $8$ to $13$; there are $15$, $136$, $741$, $3{,}200$,
$12{,}730$ and $48{,}972$ of them, $65{,}794$ in all, and none of smaller size. For each we followed the
first $24$ steps of its spine, replacing a state by a convertible reduct whose two parts are in normal
form when one was found, and recorded it as \emph{certified covered} if one of its states equals,
as written, a state of $p_K$, the start of a run covered by~\citep[Corollary~3 or~4]{s12}, or a state of
another certified term, which suffices by Remark~\ref{rem:known}. Of the $65{,}794$ terms, $535$ ($0.8\%$)
are certified covered, in six classes; for example, the seventh state of
$S\,S\,S\,(S\,(S\,S))\,(S\,(S\,S))\,S$, of size $10$, is $p_K$ as written. No term of size $8$ or $9$
is certified. The count is a lower bound: a term whose states become convertible with a covered one
only after $24$ steps, or only up to conversion, is not seen. Shared states were detected by hashing;
an independent program that compares states as written recovers the same six classes with the same
members. These counts are computations, not theorems.

Among the terms of least size is $X\,T=S\,(S\,S\,S)\,S\,(S\,S\,S)$, of size $8$. Whether $\CT(X\,T)$
holds is open. The census certifies CT at no term of size $8$ or $9$. The terms $R(n,j)$ of
Proposition~\ref{prop:afam}, of sizes $50$ and more, are much larger, and CT at them is open too.
\end{remark}

\section{Scope}\label{sec:scope}

\emph{Open problems.} Lock~1, CT and NTD are open. Theorem~\ref{thm:main} proves neither CT nor NTD from
the other, and it gives no new case of Lock~1. The known cases of CT are those of~\citep{s11,s12}
(Remark~\ref{rem:known}); the known cases of Lock~1 on classes of terms that do not normalize
include those of~\citep{s14}, whose trees have an hnf at every node of their rightmost path.

\emph{Classical steps.} Only the converse direction of Theorem~\ref{thm:main} is classical. It uses
excluded middle for $\conv$ to choose a non-convertible pair of children and a least non-convertible
pair of teeth, and its last step needs Markov's principle. Every other proof is constructive.

\emph{Inputs.} Tree equality is the relation of~\citep[Definition~1]{s13}, equality of all prefixes
$\NF_d$; we do not use its reading as equality of infinitary normal forms. Lemma~\ref{lem:replace}
and Proposition~\ref{prop:bot} use Theorem~C of~\citep{s13}. Proposition~\ref{prop:afam} uses six entries
of the block table, read by a program. Section~\ref{sec:known} rests on~\citep{s12}, which is not
mechanized, and its counts are computations.

\section*{Data availability}
\sloppy
Nothing in this paper is mechanized in Lean. The inputs it cites from~\citep{s13} are mechanized in
\texttt{lean-ski}~\citep{leanski} (Lean~4~\citep{lean4}), version~1.24.0, with only the axioms
\texttt{propext} and \texttt{Quot.sound}: tree equality is \texttt{Congruence.Eqv}; Theorem~C is
\texttt{Congruence.cong\_eqv}, with the application form \texttt{Congruence.eqv\_app}; Fact~K, used
as (K2), is \texttt{Congruence.factK}; (K1) is \texttt{Congruence.eqv\_bot} and
\texttt{Congruence.eqv\_parts}; (N2) and (N4) are \texttt{InfNF.nf\_wsteps} and
\texttt{Congruence.node\_inv}; Lock~1 is \texttt{WordProblem.Lock1}. The block table is
\texttt{blockTable} in \texttt{Ski/HnfS.lean}. The lemmas of~\citep{s11} used as (L1)--(L3) are
part of the mechanization of~\citep{s11}. The census of Remark~\ref{rem:census} and the checks of
Proposition~\ref{prop:afam} were run with Python programs that read \texttt{blockTable}; they are
available from the authors on request.

\bibliographystyle{plainnat}
\bibliography{refs}
\end{document}
